\section{Conclusions}\label{sec:conclusions}

Silica scaling limits how far geothermal brine can be cooled before
reinjection, and the rate law used to set that limit differs from
independent laboratory data by a factor of \NUM{KrMed} at the median.
Under that uncertainty, and with feedback limited to silica grab
samples and one rate measurement per day, this study reached four
conclusions.
\begin{enumerate}
\item A daily correction of the scaling constraint must compare the
measured rate with the model at the measured silica concentration.
Compared with the design-case prediction, it mistakes the optimizer's
safety margin for model error: under sixfold faster kinetics the line
exceeded the scaling limit on \CL{S}{CAnom}{v} days per 45-day run,
against \CL{S}{CAk0}{v} with the measured anchor.
\item With the right anchor a scalar bias correction is enough.
Gaussian-process and gradient-boosted corrections traced overlapping
trade-offs between power and safety and did not improve on it. A
back-off of $\kappa = 1$ held violations to \CL{S}{CAk1}{v} days at a
\CL{S}{CAk1}{up}\% power gain over fixed practice under fast kinetics,
and without systematic kinetic error the correction gained
\CL{B}{CAk0}{up}\% with no violation on any plant.
\item In plant terms the safe response to fast kinetics is a few
kelvin: the corrected policy raised the mean reinjection temperature
from \CL{S}{Open}{T} to \CL{S}{CAk1}{T}~$^{\circ}$C, and the optimized
dose stayed below 1~ppm throughout.
\item A hybrid soft sensor estimates the scaling rate from sparse daily
labels with \NUM{T1redPh}\% less error than its kinetic core, provided
its correction is earned from the plant's own labels.
\end{enumerate}
The laboratory comparison also returned two findings of geochemical
interest: the decay of a single monomeric-silica curve is second order,
against the fourth-order peak-rate law, and the 2000 solubility
correlation agrees with the 2024 endpoints at a median of
$1.05\times$. A field trial on a side stream with inspection spools is
the next step.
